\section{Agent 产品判断框架}

前面几章分别讲了体验、模型变量、容器、AgentRank 和 OS Agent。本章把这些观点整理成一个可复用框架：判断一个 Agent 应用，不应只问“模型强不强”，而要问它是否定义了新任务环境、是否创造新人群、是否让 token 更高效地产生价值、是否形成网络效应，以及是否有身份和信任机制。

\begin{knowledgebox}{课堂提示：Agent 产品不是自动化脚本}
自动化脚本完成一个任务就结束；Agent 产品要在环境里持续产生任务、反馈、调用和关系。YouWare 的野心不在于一次性生成网站，而在于让 AI coding 作品可以发布、互动、被调用，并进入新的创作网络。
\end{knowledgebox}

\noindent\begin{tabularx}{\textwidth}{p{0.20\textwidth}p{0.34\textwidth}X}
\toprule
判断维度 & 要问的问题 & YouWare 对应点 \\
\midrule
新人群 & 是否出现过去不会做这件事的人？ & Vibe Coder 不等于传统程序员。 \\
容器 & 作品是否有新的承载和分享环境？ & HTML/code 变成可访问、可交互网站。 \\
智能效率 & token 是否高效转成作品和价值？ & 从 prompt 闲聊转向生成可用产物。 \\
网络效应 & 用户、作品、Agent 是否互相增强？ & 作品分发、调用、AgentRank 形成正循环。 \\
信任治理 & 是否有身份、权限、审计和安全边界？ & Agent 部落化后需要身份证和密码锁。 \\
模型依赖 & 模型变强时，产品是否自然受益？ & 永远相信 Model 会变好，但应用做环境。 \\
\bottomrule
\end{tabularx}

\begin{importantbox}{好的 Agent 应用应该顺着模型变强}
如果模型每进步一次，产品都要重写一层脚手架，说明产品站错了位置；如果模型每进步一次，环境里的作品质量、任务范围和用户价值自然放大，说明产品处在正确变量上。
\end{importantbox}

\subsection{三个常见误区}

本节把框架反过来看。第一个误区是把 Agent 产品做成聊天框。聊天框足够通用，但环境太薄，用户不知道怎样把智能转成可分享作品。第二个误区是把 Agent 产品做成传统企业软件。企业软件收入明确，但容易变重，压缩新人群和生态网络。第三个误区是把模型能力当作应用壁垒。模型公司会继续变强，应用壁垒必须在用户环境、数据回流、网络效应和信任机制里。

\begin{warningbox}{应用壁垒不能只靠先发 UI}
今天的 UI 很容易被复制。真正难复制的是用户关系、作品网络、Agent 调用关系、身份信任和持续反馈。如果这些没有形成，先发只是曝光，不是壁垒。
\end{warningbox}

\subsection{本章小结}

Agent 产品判断框架可以压缩为一句话：看它是否创造了新的任务环境和新人群，并让模型能力在这个环境中持续增值。YouWare 的核心赌注，就是 AI coding 会从工程能力变成创作能力，而创作能力需要新的容器和网络。

